\section{Discussion}
\label{sec:discussion}

\textbf{What the oracle decides.} Arm D feeds the true dock velocity through the same law as arm B, so it isolates the plant from the estimator. If D over-swings and fails at the periods where B does, the plant's gain and lag limit the envelope and no estimator could have recovered it; if D docks where B fails, the difference is estimation; and if C approaches D across the sweep, the mechanism of Section~\ref{sec:failure} and the redesign of Section~\ref{sec:redesign} are both confirmed. [Outcome sentence to be written from Section~\ref{sec:results} once the core campaign has run.]

\textbf{Transfer to hardware.} The cliff sits where the plant lag becomes comparable to the sway period, so whether it moves on a real vehicle depends on whether the identified lag is real. The autopilot in the loop is the production ArduSub firmware, so the path from command through mode controller to motor mixing is the deployed code, and the rest of the lag is the physics of accelerating the vehicle's mass and added mass under bounded thrust, taken from the vehicle's hydrodynamic model. A real vehicle adds ESC and thruster response, the autopilot's own sensor filtering and tether drag, none of which the simulation models, so the lag can only grow and the instability of an open-loop feedforward would appear at the same or longer periods. The redesign does not assume the lag away; it identifies it and places the velocity loop's bandwidth under it.

\textbf{Navigation error.} The observability limit of Section~\ref{sec:navobs} has a plain practical reading. Any onboard dock-velocity estimate carries the vehicle's own slow velocity error, whatever frame it is estimated in, so an open-loop feedforward on a real navigation solution drives the vehicle at the dock velocity plus the bias. Closing the feedforward on the same navigation solution cancels the bias, which makes the velocity loop of Eq.~\eqref{eq:vloop} the part of the redesign that matters most for a deployment, ahead of the choice of estimator.

\textbf{Limitations.} Contact is not modelled; success is scored at the capture envelope and would-be structural contacts are censored, which a funnel dock resolves mechanically but which the simulation cannot show. The dock motion is a single prescribed sinusoid at one amplitude, a reproducible benchmark rather than a sea state; the estimators are deliberately model-free about the motion so that the conclusions do not depend on it. The nominal sweep uses ground-truth navigation, which is what made the mechanism identifiable and is also why the navigation-error sweep is reported separately. The plant gain and lag quoted in Section~\ref{sec:failure} are closed-loop fits, replaced by the step-test identification for the redesign. Finally, every diagnostic the system exposes pointed at the estimator, and only the recorded ground truth exposed the plant; a deployment of any velocity feedforward needs an independent check of the vehicle's actual response, not just of the estimate it is following.
